% ============================================================
%  Поглавље 5: Основе квантног рачунарства
% ============================================================
\section{Основе квантног рачунарства}
\label{sec:kvant}

% --- Циљ поглавља ---
% Дати минималну али самодовољну основу квантног рачунарства
% неопходну за разумевање QAOA у следећем поглављу.
% Избегавати претерану физичку формалност; фокус на рачунарски аспект.

\subsection{Кубити и суперпозиција}

% --- Кубит ---
% Кубит је квантна аналогија класичног бита. Стање кубита је вектор
% у дводимензионалном комплексном Хилбертовом простору:
%   |ψ⟩ = α|0⟩ + β|1⟩,  |α|² + |β|² = 1.
% Мерење даје 0 са вероватноћом |α|² и 1 са вероватноћом |β|².

Стање кубита описује се вектором у Хилбертовом простору $\mathcal{H} = \mathbb{C}^2$:
\[
  |\psi\rangle = \alpha|0\rangle + \beta|1\rangle, \quad
  \alpha, \beta \in \mathbb{C}, \quad |\alpha|^2 + |\beta|^2 = 1.
\]

% TODO: Илустровати Блохову сферу (TikZ или слика).
% Поменути да систем од n кубита живи у простору 2ⁿ димензија.

\subsection{Квантна заплетеност}

% --- Entanglement ---
% Два кубита су заплетена ако се стање система не може записати
% као тензорски производ стања појединачних кубита.
% Пример: Белово стање (|00⟩ + |11⟩)/√2.
% Заплетеност омогућује корелације без класичне комуникације.

\begin{example}[Белово стање]
\[
  |\Phi^+\rangle = \frac{|00\rangle + |11\rangle}{\sqrt{2}}
\]
Мерење prvog кубита тренутно одређује стање другог, без обзира на растојање.
\end{example}

\subsection{Квантна кола и капије}

% --- Квантне капије ---
% Квантне капије су унитарне трансформације. Набројите:
% H (Хадамар), X (NOT), CNOT, Rz(θ), CZ.
% Нагласите да је Хадамар капија кључна за стварање суперпозиције.

\[
  H = \frac{1}{\sqrt{2}}\begin{pmatrix}1 & 1\\ 1 & -1\end{pmatrix}, \qquad
  X = \begin{pmatrix}0 & 1\\ 1 & 0\end{pmatrix}, \qquad
  R_Z(\theta) = \begin{pmatrix}e^{-i\theta/2} & 0\\ 0 & e^{i\theta/2}\end{pmatrix}.
\]

% TODO: Приказати једноставно квантно коло (нпр. Bell state) помоћу
% qcircuit пакета или TikZ.

\subsection{Квантна предност и NISQ ера}

% --- NISQ ---
% Прескил 2018: тренутни квантни рачунари имају 50--1000 зашумљених
% кубита (Noisy Intermediate-Scale Quantum, NISQ).
% Нису довољно велики за потпуну квантну корекцију грешака,
% али омогућују хибридне квантно-класичне алгоритме.

% TODO: Поменути тренутне хардверске платформе:
% IBM Quantum, Google Sycamore, IonQ, Rigetti.
% Навести тренутне рекорде у броју кубита и дубини кола.

\begin{remark}
Квантна предност за НП-тешке проблеме теоријски није доказана.
Истраживања показују да NISQ алгоритми могу пружити практичну предност
на специфичним инстанцама пре него теоријску асимптотску супериорност.
\end{remark}
